% ============================================================================
% Fichier  : partie4/P04_C13_forensique_memoire.tex
% Titre    : Forensique Mémoire
% Auteur   : MINKA MI NGUIDJOÏ Thierry Emmanuel
% ============================================================================
% Sources : NIST SP 800-86 ; RFC 3227 ; MITRE ATT&CK. Opération FARO.
% ============================================================================

\chapter{Forensique Mémoire}
\label{chap:memoire}
\chaptermark{Forensique mémoire}

\epigraph{%
    « La mémoire vive est le seul témoin qui dit la vérité du moment présent\\
    et meurt dès qu'on éteint. »%
}{Principe du cours}

\niveauChapitre{Expert}{Chapitre~\ref{chap:systeme}.}

\begin{boxobjectifs}
\begin{itemize}
    \item Comprendre pourquoi et quand capturer la mémoire vive.
    \item Identifier ce que la mémoire révèle et que le disque ignore.
    \item Décrire l'acquisition et ses précautions.
    \item Repérer les traces d'exécution en mémoire.
    \item Mesurer le coût d'un redémarrage imprudent.
\end{itemize}
\end{boxobjectifs}

\bigskip

% ============================================================================
\section*{Question directrice : qu'a-t-on perdu en redémarrant FS-01~?}
% ============================================================================

En redémarrant le serveur, le technicien a effacé la mémoire vive. Ce
chapitre montre ce qu'elle contenait, pourquoi elle prime dans l'ordre de
volatilité, et comment on la capture quand elle est encore là.

\medskip
\textbf{Plan du chapitre.}
\begin{enumerate}
    \item Pourquoi la mémoire (section~\ref{sec:c13-pourquoi}).
    \item Ce qu'elle révèle (section~\ref{sec:c13-revele}).
    \item L'acquisition (section~\ref{sec:c13-acquisition}).
    \item Le coût du redémarrage (section~\ref{sec:c13-cout}).
\end{enumerate}

\separateur

% ============================================================================
\section{Pourquoi la Mémoire}
\label{sec:c13-pourquoi}
% ============================================================================

La mémoire vive contient l'état réel du système au moment présent~: processus
en cours, connexions réseau ouvertes, clés et secrets déchiffrés, code
malveillant qui ne touche jamais le disque. C'est pourquoi elle vient en tête
de l'ordre de volatilité (Chapitre~\ref{chap:fondements})~: elle disparaît à
l'extinction.

% ============================================================================
\section{Ce qu'elle Révèle}
\label{sec:c13-revele}
% ============================================================================

\begin{tableaucours}{L{4.0cm}L{8.2cm}}{%
    \tableauHeader{Élément en mémoire} &
    \tableauHeader{Intérêt pour l'enquête}
}
Processus en cours & Code malveillant actif, parent suspect \\
Connexions réseau & Vers où le système parlait (exfiltration, pilotage) \\
Secrets déchiffrés & Clés, mots de passe, jetons en clair \\
Code injecté & Programme présent en mémoire mais absent du disque \\
Commandes récentes & Ce que l'attaquant faisait à l'instant \\
\end{tableaucours}
\vspace{-0.8em}
\begin{center}{\small\itshape Tableau~13.1. Ce que révèle la mémoire vive}\end{center}

% ============================================================================
\section{L'Acquisition}
\label{sec:c13-acquisition}
% ============================================================================

\begin{boxprocedure}
\textbf{Acquérir la mémoire vive}
\begin{enumerate}
    \item Agir sur le système vivant, avant toute extinction.
    \item Utiliser un outil d'acquisition à empreinte minimale, documenté.
    \item Écrire la capture sur un support externe, jamais sur le disque
          analysé.
    \item Calculer le condensat de la capture.
    \item Consigner l'acte dans la chaîne de custody.
\end{enumerate}
\end{boxprocedure}

\begin{boxCRO}
\textbf{Capturer la mémoire, un arbitrage CRO.} L'acquisition touche
nécessairement le système vivant~: elle laisse une empreinte \cro{0}{-}{+}.
On accepte une légère altération (fiabilité) pour sauver une trace unique et
opposable (opposabilité), et on la documente pour la rendre défendable.
\end{boxCRO}

% ============================================================================
\section{Le Coût du Redémarrage}
\label{sec:c13-cout}
% ============================================================================

\begin{boxscenario}{Opération FARO. Ce que la mémoire aurait dit}
\textbf{Constatation.} La mémoire de \texttt{FS-01}, perdue au redémarrage,
aurait vraisemblablement contenu~: le processus de chiffrement actif, la
connexion d'exfiltration ouverte, et peut-être la clé de chiffrement en clair.

\medskip
\textbf{Analyse} (confiance moyenne). Sans la mémoire, l'enquête a dû
reconstituer ces éléments indirectement, à partir du disque et du réseau~:
plus long, moins complet. La clé en clair, si elle avait été capturée, aurait
pu éviter de payer la rançon. Le redémarrage a coûté cher.
\end{boxscenario}

\begin{boxmetier}{Analyste forensique et analyste SOC}
\textbf{Capturer la mémoire d'abord.} Face à un système compromis encore
allumé, le réflexe est de capturer la mémoire avant tout autre geste. Éteindre
« pour arrêter l'attaque » peut détruire la preuve qui aurait permis de la
comprendre.
\end{boxmetier}

\separateur

\begin{boxresume}
\textbf{Ce qu'il faut retenir de ce chapitre~:}
\begin{itemize}
    \item \textbf{La mémoire vive} contient l'état réel, et meurt à
          l'extinction.
    \item \textbf{Elle révèle} processus, connexions, secrets, code injecté.
    \item \textbf{On la capture} sur le système vivant, vers un support
          externe.
    \item \textbf{L'acquisition} est un arbitrage CRO documenté.
    \item \textbf{Un redémarrage} détruit cette preuve unique.
\end{itemize}
\end{boxresume}

\begin{boxinfo}
\textbf{Points de contrôle de l'élève.} Je sais~: expliquer la primauté de la
mémoire~; citer ce qu'elle révèle~; décrire son acquisition~; analyser le CRO
de l'acquisition~; mesurer le coût d'un redémarrage.
\end{boxinfo}

\bigskip

\begin{boxexo}[title={\ding{45}\quad Exercice~13.1. Pourquoi d'abord \niveaubase}]
Pourquoi la mémoire vient-elle en tête de l'ordre de volatilité~?
\end{boxexo}

\begin{boxexo}[title={\ding{45}\quad Exercice~13.2. Révélations \niveaubase}]
Citez trois éléments que la mémoire révèle et que le disque ignore.
\end{boxexo}

\begin{boxexo}[title={\ding{45}\quad Exercice~13.3. Support \niveaubase}]
Sur quel support écrit-on la capture mémoire~? Pourquoi pas le disque
analysé~?
\end{boxexo}

\begin{boxexo}[title={\ding{45}\quad Exercice~13.4. CRO \niveaumoyen}]
Expliquez l'arbitrage CRO de l'acquisition mémoire.
\end{boxexo}

\begin{boxexo}[title={\ding{45}\quad Exercice~13.5. FARO \niveaumoyen}]
Qu'aurait pu contenir la mémoire de \texttt{FS-01}~? Qu'a coûté sa perte~?
\end{boxexo}

\begin{boxexo}[title={\ding{45}\quad Exercice~13.6. Code injecté \niveaumoyen}]
Pourquoi un code injecté peut-il être invisible sur le disque~?
\end{boxexo}

\begin{boxexo}[title={\ding{45}\quad Exercice~13.7. Éteindre \niveaumoyen}]
« Il faut éteindre pour arrêter l'attaque. » Discutez.
\end{boxexo}

\begin{boxexo}[title={\ding{45}\quad Exercice~13.8. Empreinte \niveaumoyen}]
Pourquoi documente-t-on l'empreinte laissée par l'outil d'acquisition~?
\end{boxexo}

\begin{boxexo}[title={\ding{45}\quad Exercice~13.9. Procédure \niveauexpert}]
Rédigez la procédure de capture mémoire de première intervention pour LiZbiZ.
\end{boxexo}

\begin{boxexo}[title={\ding{45}\quad Exercice~13.10. Consigne \niveauexpert}]
Rédigez la consigne « système compromis allumé~: que faire en premier~? ».
\end{boxexo}

\bigskip

\begin{boxinfo}
\textbf{Pour aller plus loin.}
\begin{itemize}
    \item NIST SP~800-86~\cite{nist80086}~; RFC~3227~\cite{rfc3227}.
\end{itemize}
\end{boxinfo}

\begin{boxinfo}
\textbf{Transition.} Reste la troisième scène~: le réseau, qui garde la trace
des échanges. Le Chapitre~\ref{chap:reseau} traite du forensique réseau.
\end{boxinfo}
